\documentclass[UTF8,11pt,a4paper]{ctexart}

\usepackage[margin=1.85cm]{geometry}
\usepackage{amsmath,amssymb,mathtools}
\usepackage{booktabs,array,longtable}
\usepackage{enumitem}
\usepackage{fancyhdr}
\usepackage{hyperref}

\setlength{\parindent}{0pt}
\setlength{\parskip}{0.45em}
\setlength{\headheight}{14pt}
\setlist[enumerate]{leftmargin=1.8em,itemsep=0.28em,topsep=0.25em}
\renewcommand{\arraystretch}{1.2}

\pagestyle{fancy}
\fancyhf{}
\lhead{Artificial Intelligence(H)}
\rhead{Mock Final Exam}
\cfoot{\thepage}

\newcommand{\pts}[1]{\hfill{\bfseries [#1 pts]}}
\newcommand{\blank}{\vspace{1.0em}}

\title{\textbf{Artificial Intelligence(H)}\\
\large Mock Final Exam}
\author{Southern University of Science and Technology}
\date{Spring 2026}

\begin{document}
\maketitle

\begin{center}
\begin{tabular}{>{\bfseries}r p{0.68\textwidth}}
Time: & 120 minutes \\
Total: & 100 points \\
Instructions: & Show necessary derivations. For search and pruning questions, follow the specified order.
\end{tabular}
\end{center}

\vspace{0.5em}
This mock paper follows the style of the reconstructed final exam. All graphs, trees, and datasets are provided in text form so that the paper is self-contained.

\section*{Problem 1: Search \pts{15}}

Consider the following undirected graph. Every edge has uniform cost $c>0$.

\[
\begin{array}{c|l}
\text{Node} & \text{Neighbors in alphabetical order}\\
\hline
A & B,C,D\\
B & A,E,F\\
C & A,D,F,G\\
D & A,C,H\\
E & B,I\\
F & B,C,G,I\\
G & C,F,H,I\\
H & D,G\\
I & E,F,G
\end{array}
\]

The start node is $A$ and the goal node is $G$. If there are multiple nodes to expand, always use alphabetical order.

\begin{enumerate}[label=(\alph*)]
\item Draw the breadth-first search tree until the goal is found. State the expansion order. \pts{5}
\item Now let $c=1$ and use the following heuristic function:
\[
\begin{array}{c|ccccccccc}
n&A&B&C&D&E&F&G&H&I\\
\hline
h(n)&4&3&2&3&3&1&0&2&1
\end{array}
\]
Draw the A* search tree until $G$ is selected for expansion. For each generated node, write $g(n)$, $h(n)$, and $f(n)=g(n)+h(n)$. \pts{7}
\item For the same heuristic table, find the range of $c$ that guarantees admissibility. Show your calculation. \pts{3}
\end{enumerate}

\section*{Problem 2: Minimax and Alpha-Beta Pruning \pts{13}}

Consider the following game tree. The root is a MAX node. Children are visited from left to right.

\[
\begin{array}{c|c|c}
\text{Root child} & \text{Second-level child} & \text{Leaf utilities}\\
\hline
A\ (\text{MIN}) & A_1\ (\text{MAX}) & 3,\ 12\\
                & A_2\ (\text{MAX}) & 8,\ 2\\
                & A_3\ (\text{MAX}) & 4,\ 6\\
\hline
B\ (\text{MIN}) & B_1\ (\text{MAX}) & 14,\ 5\\
                & B_2\ (\text{MAX}) & 2,\ 1\\
                & B_3\ (\text{MAX}) & 7,\ 3\\
\hline
C\ (\text{MIN}) & C_1\ (\text{MAX}) & 6,\ 9\\
                & C_2\ (\text{MAX}) & 4,\ 0\\
                & C_3\ (\text{MAX}) & 10,\ 11
\end{array}
\]

\begin{enumerate}[label=(\alph*)]
\item Compute the minimax value of every internal node. \pts{4}
\item Run alpha-beta pruning and write the final $\alpha$ and $\beta$ values for each visited internal node. If a node is pruned, write NA. \pts{5}
\item Mark which branches are pruned. \pts{2}
\item Mark the final path chosen by MAX. \pts{2}
\end{enumerate}

\section*{Problem 3: CSP and AC-3 \pts{12}}

\begin{enumerate}[label=(\alph*)]
\item State the worst-case time complexity of AC-3 and briefly explain why. \pts{4}
\item Consider a CSP with variables $X,Y,Z$:
\[
D_X=D_Y=D_Z=\{1,2,3,4\},\qquad C=\{X<Y,\ Y<Z\}.
\]
Use AC-3 to enforce arc consistency. Show each domain pruning step and the final domains. \pts{8}
\end{enumerate}

\section*{Problem 4: Logic and CNF \pts{12}}

\begin{enumerate}[label=(\alph*)]
\item Convert the following propositional formula into CNF:
\[
(p\land q)\to \neg(p\leftrightarrow r).
\]
Simplify the result as much as possible. \pts{5}
\item Four robots $R_1,R_2,R_3,R_4$ may be selected for a mission. Let $x_i$ mean robot $R_i$ is selected. Represent the following rules in propositional logic and convert them into CNF. \pts{5}
\begin{enumerate}[label=(\roman*)]
\item $R_1$ and $R_3$ must be selected together.
\item Exactly one of $R_2$ and $R_4$ is selected.
\item If $R_2$ is selected, then $R_1$ must also be selected.
\end{enumerate}
\item Suppose additionally $x_2$ is true. Use unit propagation to derive which of $x_1,x_3,x_4$ must be true or false. \pts{2}
\end{enumerate}

\section*{Problem 5: Perceptron and Logistic Regression \pts{10}}

\begin{enumerate}[label=(\alph*)]
\item Write the perceptron prediction rule and weight update rule for labels $y\in\{-1,+1\}$. \pts{3}
\item Given $w=(-1,1)^T$, $b=0$, learning rate $\eta=0.5$, and one training example $x=(2,1)^T$, $y=+1$, perform one perceptron update if needed. \pts{3}
\item Derive the gradient descent update for binary logistic regression with cross-entropy loss:
\[
J(w,b)=-\frac1m\sum_{i=1}^m\left[y_i\log h_i+(1-y_i)\log(1-h_i)\right],
\quad h_i=\sigma(w^Tx_i+b).
\]
\pts{4}
\end{enumerate}

\section*{Problem 6: SVM \pts{10}}

Consider two one-dimensional data points $(x_1,y_1)=(2,+1)$ and $(x_2,y_2)=(-2,-1)$.

\begin{enumerate}[label=(\alph*)]
\item Write the primal objective function and constraints of the soft-margin SVM. Use $C$ as the penalty parameter. \pts{4}
\item For these two points, write the constraints explicitly in terms of $w,b,\xi_1,\xi_2$. \pts{2}
\item Suppose the optimal $w$ is given by $w^*=\min(C,\frac12)$. Find the possible optimal value(s) of $b$ for the cases $C\ge\frac12$ and $0<C<\frac12$. \pts{4}
\end{enumerate}

\section*{Problem 7: Naive Bayes \pts{12}}

The following dataset has four binary features and one binary label.

\[
\begin{array}{c|cccc|c}
\text{ID} & X_1&X_2&X_3&X_4&Y\\
\hline
1&1&0&1&0&+\\
2&1&1&1&0&+\\
3&0&1&0&1&-\\
4&0&0&0&1&-\\
5&1&0&0&0&+\\
6&0&1&1&1&-\\
7&1&1&0&0&+\\
8&0&0&1&1&-
\end{array}
\]

Use Bernoulli Naive Bayes with Laplace smoothing $\alpha=1$.

\begin{enumerate}[label=(\alph*)]
\item Compute the class priors $P(Y=+)$ and $P(Y=-)$. \pts{2}
\item Compute all conditional probabilities needed for prediction, i.e. $P(X_j=1\mid Y)$ and $P(X_j=0\mid Y)$ for all $j$. \pts{5}
\item Predict the label of $x=(1,0,1,0)$. Show both class scores. \pts{3}
\item If the Naive Bayes independence assumption is not used, how many feature-label combinations must be considered for four binary features and one binary label? How many independent parameters does the full joint distribution over the five binary variables have? \pts{2}
\end{enumerate}

\section*{Problem 8: Bayesian Networks and Sampling \pts{14}}

Suppose $A\to B\to C\to D$ is a Bayesian network with:
\[
P(A)=\frac14,\quad
P(B\mid A)=\frac15,\quad P(B\mid \neg A)=\frac34,
\]
\[
P(C\mid B)=\frac12,\quad P(C\mid \neg B)=\frac13,\quad
P(D\mid C)=\frac16,\quad P(D\mid \neg C)=\frac78.
\]

\begin{enumerate}[label=(\alph*)]
\item Factorize the joint distribution $P(A,B,C,D)$ according to the network. \pts{2}
\item Given these prior samples:
\[
\begin{aligned}
&(A,B,\neg C,\neg D),\quad (A,\neg B,C,\neg D),\\
&(\neg A,B,C,\neg D),\quad (\neg A,\neg B,C,\neg D),\\
&(A,\neg B,\neg C,D),\quad (A,B,C,\neg D),\\
&(\neg A,B,\neg C,D),\quad (\neg A,\neg B,C,\neg D),
\end{aligned}
\]
estimate $P(C)$ by prior sampling. \pts{2}
\item Use rejection sampling on the same samples to estimate $P(C\mid A,\neg D)$. State which samples are kept. \pts{3}
\item For likelihood weighting to estimate $P(\neg A\mid B,\neg D)$, compute the weight of each sample:
\[
\neg A,B,C,\neg D;\quad A,B,C,\neg D;\quad A,B,\neg C,\neg D;\quad \neg A,B,\neg C,\neg D.
\]
\pts{4}
\item From the four weighted samples in part (d), estimate $P(\neg A\mid B,\neg D)$. \pts{3}
\end{enumerate}

\section*{Problem 9: Genetic Algorithm and Course Projects \pts{12}}

\begin{enumerate}[label=(\alph*)]
\item Suppose we evolve the weights of a neural network using a genetic algorithm instead of backpropagation. Describe how to set up a fair experiment comparing these two training methods. \pts{5}
\item In Project 1, the objective of information exposure maximization can be estimated by Monte Carlo simulation. Explain what the state/individual, fitness function, and stopping criterion could be for a GA-based solver. \pts{4}
\item In Project 3, a recommender system is evaluated by AUC and nDCG@5. Briefly explain what each metric measures and why a model may improve one but not the other. \pts{3}
\end{enumerate}

\end{document}
